\section{Introduction}

Selling the last deed in a rival's color group raises cash now and may create a larger rent threat later. Buying it can leave too little money to build. Keeping its owner solvent can preserve a blocker against a third player; refusing to finance that owner can deliver the same property through bankruptcy. These decisions require an agent to connect a proposed exchange to other players' opportunities and to its own future actions.

LLM negotiation research has studied agreement, utility and conversational tactics in settings with explicit payoffs~\cite{lewis2017deal,he2018decoupling,bianchi2024negotiationarena}. Complete economic games add persistent consequences: an agreement changes the next negotiation, construction changes future debts, and elimination changes who owns the productive assets. The question is therefore not only whether an agent can negotiate plausible terms, but which consequences it attends to and whether it can carry a financial plan into action. Strategic attention and strategic effectiveness are different objects. An agent can correctly identify a rival's opportunity without choosing a good price or a good refusal.

We examine this distinction in twenty complete MonopolyBench games played by GPT-6.1 Sol, Gemini 4 Argon, Claude Opus 5.5 and Grok 4.7. The corpus contains 6,989 decisions, 526 negotiation threads and 188 cash-shortfall episodes requiring explicit liquidation. Every offer is linked to the recipient's delivered state and response, and every financing sequence to the resulting payment or bankruptcy. This supports an observational account of bargaining and economic execution, rather than a leaderboard of underlying model ability.

The central bargaining pattern depends on what the proposer offers in return. Opening swaps completing only the proposer's group receive 81.9\% rejections, compared with 51.4\% for swaps completing neither party. Cash bids for completing deeds do not receive more rejections than other cash bids. Asking prices nevertheless vary with the buyer's holdings, and statements accompanying both refusals and prices discuss construction budgets, financing and prospective estates. The agents attend to public economic consequences, but their behavior is not captured by a single rule of withholding every completing deed. The holder's cash and the negotiation history further qualify the pattern.

Financial execution supplies a complementary test with a known endpoint. Agents pay 128 of 131 debts for which payment is demonstrated by actual play or earlier legal-sequence checks. In 19 successful episodes, selling buildings makes new mortgages available and the agent uses that collateral. Two payable bankruptcies accompany calculations that count the initial mortgage menu and building salvage but omit the mortgages those sales would unlock. The contrast concerns a fundamental planning problem: anticipating how an action changes the next available actions.

The paper contributes three linked empirical results:
\begin{enumerate}
\item A complete negotiation account showing how the completion contrast differs between swaps and cash bids, with opening stage, holder cash and recipient configuration made explicit.
\item A census of financing sequences that distinguishes successful local execution, unpayable obligations, and failures to anticipate newly available collateral.
\item Eight contextual incidents connecting refusal, pricing, bluffing, factual correction, selective disclosure, a financing-denial appeal and withdrawal of an assurance to subsequent action, including a completed ten-action payment plan.
\end{enumerate}

The incidents sit within the economic results they illuminate. Their later consequences matter: the player who buys or builds a group may eventually transfer it to the rival it sought to contain. They are selected examples, not estimates of deception frequency. Public deception is permitted by the task; calling a tactic a bluff does not make it an instruction violation.

The broader research agenda asks about strategic reasoning, opponent modeling, negotiation and deception, planning and memory, adaptation, risk management, collective dynamics, and evaluation validity. It also asks whether strategic talk becomes action, whether general capability predicts multi-agent capability, and how skill can be separated from chance. Appendix~\ref{app:questions} gives a direct answer at the level supported for every question. The present study is strongest on observable bargaining and financial execution. It has no independent general-capability measure, intervention on memory or communication, or value reference for rejected trades.
